% !PRINTED 82
\weberpage{82}
%<<CONT

\begin{equation}
\begin{aligned}
 a_1^{(1)}x_1+a_2^{(1)}x_2+\cdots+a_n^{(1)}x_n&=0,\\
 a_1^{(2)}x_1+a_2^{(2)}x_2+\cdots+a_n^{(2)}x_n&=0,\\
 &\ \vdots\\
 a_1^{(m)}x_1+a_2^{(m)}x_2+\cdots+a_n^{(m)}x_n&=0.
\end{aligned}
\tag{1}
\end{equation}
where the coefficients \(a_i^{(k)}\) are regarded as given quantities. For the present we shall make no assumption whatever about the numbers \(m,n\), but set ourselves the general task of determining \gesperrt{all systems of values} of \(x_1,x_2,\ldots,x_n\) which satisfy equations (1).

We can give one solution of equations (1) at once: whatever the coefficients \(a_i^{(k)}\) may be, they are satisfied when
\begin{equation}
        x_1=0,\quad x_2=0,\quad\ldots,\quad x_n=0.
\tag{2}
\end{equation}
We may also mention another extreme case: if all the coefficients \(a_i^{(k)}\) have the value zero, then equations (1) are satisfied for arbitrary values of \(x_1,x_2,\ldots,x_n\).

We precede the general answer to the question by the following remarks. We write the system of coefficients of (1) in the form of a rectangle
\begin{equation}
\begin{array}{cccc}
 a_1^{(1)}&a_2^{(1)}&\cdots&a_n^{(1)}\\
 a_1^{(2)}&a_2^{(2)}&\cdots&a_n^{(2)}\\
 \vdots&\vdots&&\vdots\\
 a_1^{(m)}&a_2^{(m)}&\cdots&a_n^{(m)}.
\end{array}
\tag{3}
\end{equation}
Such a scheme, which by itself has no numerical meaning, is called a \gesperrt{matrix}, insofar as it is regarded as the source of a larger number of determinants. \gesperrt{The determinants arising from the matrix are obtained by deleting arbitrary rows and columns, in an arbitrary number limited only so far that the remaining elements form a square, and by regarding this square as a determinant.} Thus one obtains from the matrix one-rowed, two-rowed, and so forth, determinants. The highest determinants are \(n\)- or \(m\)-rowed, according as \(n\) or \(m\) is the smaller number (or \(n\)-rowed if \(n=m\)).
